%% ma: L10-B14-0028
%% lop: 10
%% chuong: 5
%% bai: 14
%% dang: 10.14.5
%% loai: DS
%% muc: [NB, TH, TH, TH]
%% dap_an: "ĐSSS"
%% nguon: "chuyen-de-cac-so-dac-trung-cua-mau-so-lieu-khong-ghep-nhom-toan-10-chuong-trinh-moi.pdf (GV. Phan Nhật Linh) (P04), Bài 03 Dạng 1 Phần II Câu 5"
%% kien_thuc: "Lớp 10 Bài 13 (số trung bình, tứ phân vị), Bài 14 (khoảng biến thiên, khoảng tứ phân vị, giá trị bất thường)"
%% trang_thai: chinh-thuc
%% ghi_chu: "Tự duyệt theo yêu cầu giáo viên 22/10/2026."
\begin{ex}
Mẫu số liệu sau cho biết chiều cao (đơn vị: cm) của các bạn học sinh tổ 1 lớp 10A:
\begin{center}
\begin{tabular}{|c|c|c|c|c|c|c|c|}\hline
$163$ & $159$ & $172$ & $167$ & $165$ & $168$ & $170$ & $161$\\\hline
\end{tabular}
\end{center}
Xét tính đúng sai của các khẳng định sau:
\choiceTF{\True Khoảng biến thiên của mẫu số liệu là $13$}{Chiều cao trung bình của các học sinh tổ 1 là $160$ cm}{Khoảng tứ phân vị của mẫu số liệu là $\Delta_Q=10$}{Nếu cộng mỗi giá trị của mẫu số liệu với $2$ thì khoảng biến thiên tăng lên $2$ đơn vị}
\loigiai{Sắp xếp: $159$; $161$; $163$; $165$; $167$; $168$; $170$; $172$.

a) Đúng: $R=172-159=13$.

b) Sai: $\overline x=\dfrac{1325}8=165{,}625$ (cm).

c) Sai: $Q_1=\dfrac{161+163}2=162$; $Q_3=\dfrac{168+170}2=169$; $\Delta_Q=7$.

d) Sai: $R=(172+2)-(159+2)=13$, không đổi.}
\end{ex}
